\section{Il citoscheletro}

% ---------------------------------------------------------------------------
\subsection{Struttura e funzioni}

Il \textbf{citoscheletro} è un sistema di filamenti proteici; è una struttura \textbf{dinamica}, continuamente scomposta e riassemblata.

Funzioni:
\begin{itemize}
  \item mantiene la struttura e la forma della cellula (\textbf{struttura e supporto});
  \item fornisce ancoraggio per organuli e macromolecole (\textbf{organizzazione spaziale});
  \item permette il \textbf{trasporto intracellulare};
  \item nella divisione cellulare forma le fibre del fuso mitotico;
  \item permette il movimento cellulare (\textbf{contrattilità e motilità}).
\end{itemize}

Tre componenti (reti di fibre di diverso tipo):
\begin{itemize}
  \item \textbf{microtubuli}: $\phi$ 20--25 nm;
  \item \textbf{filamenti intermedi}: $\phi$ 8--12 nm;
  \item \textbf{microfilamenti}: $\phi$ 5--7 nm.
\end{itemize}

% ---------------------------------------------------------------------------
\subsection{Microtubuli}

\begin{itemize}
  \item diametro 20--25 nm; costituiti da subunità proteiche di \textbf{tubulina} ($\alpha$ e $\beta$);
  \item agiscono da impalcatura per determinare la forma della cellula;
  \item permettono il movimento degli organelli e delle vescicole;
  \item formano le fibre del fuso per separare i cromosomi durante la mitosi;
  \item all'interno di ciglia e flagelli permettono il movimento cellulare.
\end{itemize}

\rubrica{Formazione e polarità}
\begin{itemize}
  \item si formano per aggiunta di eterodimeri di $\alpha$-tubulina e $\beta$-tubulina a un'estremità del cilindro cavo;
  \item il cilindro possiede una \textbf{polarità}: a un'estremità (estremità $+$) si aggiungono dimeri, all'estremità opposta (estremità $-$) si rimuovono;
  \item per ogni giro di spirale sono necessari 13 dimeri.
\end{itemize}

\rubrica{Proteine associate ai microtubuli (MAP)}
\begin{itemize}
  \item presentano 3 siti di legame alla tubulina, intervallati da una distanza sufficiente a permettere il legame a 3 subunità separate;
  \item le code delle MAP si proiettano all'esterno, dove possono interagire con altri componenti cellulari.
\end{itemize}

\rubrica{Proteine motrici: la chinesina}
Proteina motrice dei microtubuli.
\begin{itemize}
  \item struttura ($\sim$80 nm): 2 catene pesanti avvolte a spirale nella regione a stelo e 2 catene leggere legate alla parte terminale delle catene pesanti; regioni: teste (sito catalitico), collo, stelo, coda;
  \item meccanismo: la molecola ``cammina'' lungo il microtubulo alternando l'attacco e il rilascio del ``piede''; le due teste effettuano movimenti identici ma alternati (passo di 16 nm).
\end{itemize}

\rubrica{Proteine motrici: la dineina citoplasmatica}
Contiene due catene pesanti e alcune catene leggere e intermedie; ciascuna catena pesante ha una grande testa globulare generante forza, un peduncolo con sito di legame per il microtubulo e uno stelo.
\begin{itemize}
  \item lungo lo stesso microtubulo le due proteine muovono le vescicole in direzioni opposte: la \textbf{chinesina} verso l'estremità $+$ (movimento anterogrado), la \textbf{dineina} verso l'estremità $-$ (movimento retrogrado);
  \item le proteine motrici sono unite alla membrana della vescicola tramite un intermediario: la chinesina tramite proteine di membrana (integrali o periferiche), la dineina tramite un complesso di proteine solubili detto \textbf{dinactina};
  \item mediano il trasporto di vescicole, gruppi vescicolari-tubulari (VTC) e organelli.
\end{itemize}

% ---------------------------------------------------------------------------
\subsection{Filamenti intermedi}

\begin{itemize}
  \item provvedono alla resistenza meccanica delle cellule soggette a stress fisici (cellule muscolari, epiteliali, neuroni);
  \item sono spesso interconnessi ad altri filamenti del citoscheletro da sottili ponti trasversali proteici (es.\ la plectina);
  \item bastoncini flessibili di $\sim$10 nm di diametro, costituiti da \textit{protofilamenti}, a loro volta formati da subunità proteiche avvolte a elica.
\end{itemize}

\rubrica{Assemblaggio}
\begin{enumerate}
  \item ciascun monomero è costituito da domini terminali globulari separati da una lunga regione ad $\alpha$-elica;
  \item coppie di monomeri si associano parallelamente, con le estremità allineate, a formare dimeri (omodimeri o eterodimeri);
  \item i dimeri si associano in maniera antiparallela e sfalsata a formare tetrameri;
  \item i tetrameri si associano lateralmente per formare un'unità di filamento ($\sim$60 nm);
  \item i filamenti intermedi molto lunghi sono formati dall'associazione di tali unità di lunghezza.
\end{enumerate}

\rubrica{Nella cellula epiteliale}
I filamenti intermedi si irradiano per tutta la cellula, ancorati sia alla superficie esterna del nucleo sia alla superficie interna della membrana plasmatica. Le connessioni con il nucleo avvengono tramite proteine che attraversano entrambe le membrane dell'involucro nucleare; quelle con la membrana plasmatica tramite siti specializzati come i desmosomi e gli emidesmosomi.

% ---------------------------------------------------------------------------
\subsection{Microfilamenti}

\begin{itemize}
  \item diametro 5--7 nm; costituiti da due catene di molecole di actina intrecciate tra loro (doppia elica);
  \item composti dalle subunità globulari della proteina \textbf{actina};
  \item in presenza di ATP i monomeri di actina polimerizzano formando un filamento flessibile a elica.
\end{itemize}

\rubrica{Proteine che legano l'actina}
\begin{enumerate}
  \item \textbf{di nucleazione}: intervengono nel primo passo di formazione del filamento, riunendo nel giusto orientamento 2 o 3 monomeri di actina;
  \item \textbf{che sequestrano i monomeri}: impediscono la polimerizzazione dei monomeri in filamenti; regolano l'equilibrio monomero--polimero;
  \item \textbf{di incappucciamento}: regolano la lunghezza dei filamenti legandosi a una delle due estremità;
  \item \textbf{che polimerizzano i monomeri}: promuovono la crescita dei filamenti;
  \item \textbf{che depolimerizzano i filamenti}: consentono la depolimerizzazione e un rapido turnover nei siti di dinamici cambiamenti del citoscheletro (locomozione cellulare, fagocitosi, citocinesi);
  \item \textbf{che formano legami crociati}: alterano l'organizzazione tridimensionale di una popolazione di filamenti (fascicolazione);
  \item \textbf{che tagliano i filamenti}: si legano a un filamento esistente e lo rompono in due;
  \item \textbf{che legano la membrana}: si legano alla membrana plasmatica permettendone i movimenti.
\end{enumerate}

\rubrica{Microfilamenti nel microvillo}
I microvilli si trovano sulla superficie apicale di epiteli con funzione di assorbimento (rivestimento interno dell'intestino, parete dei tubuli renali). I filamenti di actina sono mantenuti in disposizione ordinata da proteine che formano fasci, come la \textbf{villina} e la \textbf{fimbrina}; la \textbf{miosina I} è presente tra la membrana plasmatica del microvillo e i filamenti periferici di actina.

% ---------------------------------------------------------------------------
\subsection{Confronto tra i tre componenti}

\begin{table}[htbp]
  \centering \footnotesize \renewcommand{\arraystretch}{1.25}
  \begin{tabularx}{\textwidth}{|l|>{\raggedright\arraybackslash}X|>{\raggedright\arraybackslash}X|>{\raggedright\arraybackslash}X|}
    \hline
     & \textbf{Microtubuli} & \textbf{Filamenti intermedi} & \textbf{Filamenti di actina} \\
    \hline
    Subunità incorporata & eterodimero di GTP--$\alpha\beta$-tubulina & circa 70 proteine diverse & monomeri di ATP-actina \\ \hline
    Sito di incorporazione & estremità $+$ ($\beta$-tubulina) & interno & estremità $+$ (sfrangiata) \\ \hline
    Polarità & sì & no & sì \\ \hline
    Attività enzimatica & GTPasi & nessuna & ATPasi \\ \hline
    Proteine motrici & chinesine, dineine & nessuna & miosine \\ \hline
    Proteine associate & MAP & plachine & proteine leganti actina \\ \hline
    Struttura & tubo rigido, cavo, inestensibile & filamenti resistenti, flessibili, estensibili & filamenti a elica, flessibili, inestensibili \\ \hline
    Dimensioni & diametro esterno 25 nm & diametro 10--12 nm & diametro 8 nm \\ \hline
    Distribuzione & tutti gli eucarioti & animali & tutti gli eucarioti \\ \hline
    Funzioni primarie & supporto, trasporto intracellulare, organizzazione della cellula & supporto strutturale & motilità, contrattilità \\ \hline
    Distribuzione sub-cellulare & citoplasma & citoplasma e nucleo & citoplasma \\ \hline
  \end{tabularx}
\end{table}

% ---------------------------------------------------------------------------
\subsection{Centrioli}

\begin{itemize}
  \item diametro $\approx$ 0,2\,$\mu$m; lunghezza $\approx$ 0,4\,$\mu$m;
  \item strutture cilindriche che contengono 9 fibrille regolarmente spaziate, ognuna formata da 3 microtubuli (arrangiamento 9$\times$3), localizzate vicino al nucleo;
  \item sono presenti quasi sempre in coppia, ad angolo retto l'uno rispetto all'altro;
  \item durante la divisione cellulare si duplicano, dando origine a 2 coppie di centrioli che si spostano ai poli opposti della cellula, dove funzionano come centri per l'organizzazione dei microtubuli che formano il fuso mitotico.
\end{itemize}

% ---------------------------------------------------------------------------
\subsection{Centrosoma}

Struttura complessa che contiene 2 centrioli circondati da materiale pericentriolare (PCM), amorfo e opaco agli elettroni, in cui avviene la nucleazione dei microtubuli.
\begin{itemize}
  \item è il principale sito di inizio dei microtubuli (MTOC) e si trova tipicamente al centro della cellula, in una fitta rete di microtubuli;
  \item si duplica una volta per ciclo cellulare, in coordinazione con la replicazione del DNA: ciascun centriolo genera un centriolo figlio, e alla mitosi le due coppie migrano ai poli opposti.
\end{itemize}

% ---------------------------------------------------------------------------
\subsection{Ciglia e flagelli}

\begin{itemize}
  \item sottili organelli motori, simili a peli, che si proiettano dalla superficie di vari tipi di cellule;
  \item sono sottili filamenti avvolti da un'estroflessione della membrana plasmatica;
  \item alla base sono ancorati dal centriolo basale (corpo basale), da cui si irradia in profondità la radice del ciglio;
  \item le \textbf{ciglia} sono più corte e più numerose rispetto ai \textbf{flagelli}.
\end{itemize}

\rubrica{Battito}
\begin{itemize}
  \item le ciglia alternano una \textit{fase attiva} (di spinta) e una \textit{fase di ritorno}; il battito, flessibile, inizia alla base e si propaga verso l'apice, con ritmo \textit{metacronale} (le ciglia di file adiacenti sono in stadi diversi del battito);
  \item meccanismo: i bracci di dineina muovono i microtubuli formando e rompendo i legami con i microtubuli adiacenti, così che ciascuno ``cammina'' lungo quello vicino; le proteine di connessione (nexina) impediscono lo scivolamento oltre un certo limite $\rightarrow$ l'azione motrice si traduce nella flessione dei microtubuli (movimento oscillatorio);
  \item i \textbf{flagelli} si muovono con una forma d'onda (asimmetrica nel caso dell'alga \textit{Chlamydomonas}); es.\ il flagello degli spermatozoi.
\end{itemize}

% ---------------------------------------------------------------------------
\subsection{Assonema}

Struttura interna di ciglia e flagelli: microtubuli raggruppati in 9 coppie (doppiette) disposte alla periferia più una coppia al centro $\rightarrow$ complesso ``9$+$2''.

\rubrica{Sezione trasversale}
\begin{itemize}
  \item 9 doppiette esterne di microtubuli, ciascuna formata da un \textbf{tubulo A} e un \textbf{tubulo B}, più una coppia centrale racchiusa nella guaina centrale;
  \item \textbf{bracci di dineina} (interno ed esterno) sul tubulo A;
  \item \textbf{raggi} (spoke) diretti verso il centro;
  \item \textbf{ponte tra le coppie} (nexina), che collega le doppiette adiacenti.
\end{itemize}

\rubrica{Sezione longitudinale}
Raggi, coppie periferiche e guaina centrale; i bracci di dineina si ripetono lungo l'asse con periodo di 96 nm (spaziature 24, 32, 40, 24 nm).

\rubrica{Corpo basale}
Alla base dell'assonema; deriva dai centrioli (arrangiamento 9$\times$3 di triplette) e ancora ciglia e flagelli alla cellula.
